\begin{abstract}
Este reporte documenta una implementación personalizada de U-Net en PyTorch para segmentación de vasos retinianos sobre imágenes de fondo de ojo. El proyecto integra tres variantes arquitectónicas desarrolladas dentro del repositorio (U-Net, Attention U-Net y U-Net++), un estudio de ablación sobre funciones de pérdida y diseño de arquitectura, una evaluación \emph{in-distribution} en DRIVE y un experimento de generalización entrenando en DRIVE y evaluando en CHASE\_DB1. Además, se analiza el efecto de una estrategia de adaptación de dominio basada en CLAHE sobre el canal verde y \emph{test-time augmentation}. Como plantilla de entrega, se incluyen resultados representativos coherentes con la configuración del proyecto: el mejor modelo alcanza un F1 de 0.828 y un AUC-ROC de 0.978 en DRIVE, mientras que en CHASE\_DB1 se observa una caída a F1=0.791. El análisis cualitativo muestra que los capilares finos concentran la mayor parte de los falsos negativos debido a su bajo contraste y a la pérdida de detalle inducida por el \emph{downsampling} del encoder.\
\end{abstract}
